\section{Related Literature} \label{sec:lit_review}

This paper sits at the intersection of three literatures: the long-running tradition on the capitalization of public capital into property values, a more recent body of work on information frictions and slow capitalization in asset markets, and a fast-growing literature on the use of computer-vision methods to construct economic measurements from high-resolution imagery. We discuss each in turn and locate our contribution.

\subsection*{Public Capital and Capitalization}

Following the cross-jurisdictional regressions of \citet{Oates1969} and the Oates Decentralization Theorem \citep{Oates1972}, the subsequent literature has been overwhelmingly concerned with the capitalization of {\it positive} shocks to the public-capital stock. \citet{cellini2010value} provide the modern reference, studying close-margin school-bond referendums in California with a dynamic regression-discontinuity design and reporting substantial capitalization of new capital projects into housing prices. \citet{biasi_what_2025} sharpen the dynamic-RD methodology and isolate the share of capitalization attributable to teacher-quality investments. \citet{edel1974taxes} is among the few early studies to attempt to estimate the role of {\it maintenance} spending using cross-jurisdictional Boston data; the cross-sectional design recovers null effects, and the dimension of public-capital depreciation goes largely unexamined in the subsequent literature.

The transportation-infrastructure literature has likewise focused on new construction. \citet{gonzalez2016paving} documents that paving previously-unpaved roads in Mexico produces a 28-percent appraisal-based increase in property values. \citet{theisen2021road} estimates a 13-percent house-price increase associated with a new highway in Norway. \citet{asher2020} documents that close-margin rural roads in India have no detectable effect on village incomes or assets, attributing prior estimates to selection on growth potential. \citet{gertler2024road} estimates effects of road maintenance in Indonesia, with mixed evidence on long-run local economic development. The classical theoretical reference for road-maintenance dynamics is \citet{rioja2003}, whose dynamic general-equilibrium calibration provides the law of motion for public-capital depreciation that we adopt in Section \ref{sec:mech}. \citet{ramey2021infrastructure} surveys the macroeconomic literature on the consequences of infrastructure investment, emphasizing that the multiplier is largest when starting capital is well below the optimal level --- a finding that resonates with our welfare estimates for the case of disinvestment from existing capital.

Two relatively recent papers approach the public-services capitalization question from angles directly relevant to our setting. \citet{brinkman2024freeway} use historical highway-revolt episodes to identify capitalization of {\it both} positive and negative shocks to neighborhood transportation infrastructure, broadening the literature beyond positive-shock identification. \citet{brasington2024fire} use a close-election RD to study capitalization of fire-protection services and explicitly document a {\it limited-awareness} mechanism in which capitalization is muted because residents are slow to register changes in public-service quality. Their evidence is the closest existing analog to the mechanism we develop, and we view our findings as complementary: in their setting, the public-service shock is binary (fire-station closure) and observable; in ours, the underlying public-capital depreciation is continuous and only gradually observable, generating a delayed-capitalization pattern absent from their setting.

\subsection*{Information Frictions and Slow Capitalization}

A separate strand of literature establishes that markets capitalize information about asset fundamentals with substantial delay when the information itself is non-salient, costly to acquire, or arriving in noisy increments. \citet{chetty2009salience} and \citet{finkelstein2009eztax} provide the canonical evidence from the tax-salience literature: agents respond more weakly to non-salient tax changes than to mechanically-equivalent salient ones, and the wedge is large enough to overturn welfare conclusions derived from fully-rational benchmarks. \citet{bakkensen2022underwater} extend this logic to housing markets, showing that under-capitalization of long-horizon flood risk is correlated with belief heterogeneity in the relevant local population. \citet{ouazad2024frictions} provide the most directly relevant theoretical apparatus, deriving conditions under which amenity capitalization is time-varying in the presence of arbitrage costs and information frictions. \citet{bishop2019hedonic} provide methodological guidance on when a dynamic hedonic specification is necessary in the presence of time-varying attributes.

Our paper contributes to this literature in three respects. First, our quasi-experimental design generates exogenous variation in a fiscal decision whose long-run consequence (physical decay of the public-capital stock) is observable to all parties only with a substantial visibility lag. This is a particularly clean laboratory for testing rational-expectations asset pricing against limited-attention or signal-extraction alternatives. Second, our novel measurement of the underlying physical state of the public-capital stock --- via computer vision on aerial imagery (Section \ref{sec:data}) --- allows us to time the observable signal precisely and demonstrate empirically that capitalization tracks the signal rather than the underlying fiscal shock. Third, the belief-updating microfoundation in Section \ref{sec:belief_updating} nests the rational-expectations benchmark as a knife-edge case and identifies the friction parameter from a single transparent moment in the data, providing a structural counterpart to the reduced-form rejection of RE.

\subsection*{Computer Vision in Economics}

Our paper contributes to a fast-growing literature on the use of computer-vision methods to construct economic measurements from high-resolution imagery. \citet{naik2017computer} use street-view imagery and a convolutional architecture to predict physical urban change, demonstrating that physical-appearance features carry forward-looking economic content. \citet{glaeser2018computervision} apply convolutional networks to real-estate imagery and document that visual features predict property values beyond conventional hedonic controls. \citet{yeh2020satellite} use Landsat imagery and deep learning to estimate economic well-being in Africa, validating the methodology against ground-truth survey data and demonstrating its scalability across geographies. \citet{currier2023} construct a Road Roughness Index from vertical-acceleration signals collected by an Uber smartphone app, providing the closest existing analog to our outcome variable but limited to areas with active Uber service.

Two prior measurement traditions exist for road quality. The first relies on engineering benchmarks such as the International Roughness Index (IRI) and the Pavement Condition Index (PCI), both of which require costly on-site inspection (profilometer vans, trained surveyor visits) and are updated irregularly with substantial geographic gaps. The second relies on smartphone-sourced acceleration data \citep{currier2023}, which has broad coverage in metropolitan areas but is sparse in suburban and rural settings precisely where many of the renewal-tax referendums in our sample occur. Our contribution is to apply a fine-tuned ConvNeXt-V2 vision model to the aerial-imagery panel from the National Agriculture Imagery Program (NAIP) and the Ohio Statewide Imagery Program (OSIP3), fine-tuned on the labeled road-surface dataset of \citet{brewer2021}, and generating consistent road-quality measurements at sub-county resolution across thirty years of imagery. To our knowledge this is the first measurement of {\it dynamic changes} in road quality at this resolution and over this horizon.

% \subsection*{Our Contribution}

% This paper contributes three pieces to the intersection of these literatures.

% First, we provide the first credibly-identified estimates of the capitalization of public-capital {\it depreciation} into property values. The vast majority of the capitalization literature is concerned with positive shocks (new bonds, new highways, school referendums); the negative-shock case has been substantially understudied despite its first-order policy importance in an aging infrastructure stock with annual depreciation of approximately 3 percent and a cumulative ASCE-estimated underinvestment of \$2.6 trillion through 2029 \citep{asce2025report}.

% Second, we develop a novel computer-vision-based measurement of road quality at sub-county resolution over a thirty-year panel, enabling year-by-year tracking of the underlying public-capital stock at a temporal frequency that existing engineering and smartphone-based measures cannot deliver outside of dense metropolitan settings.

% Third, we develop a formal test and rejection of rational-expectations capitalization in a salient local asset market and provide a tractable behavioral alternative --- a belief-updating microfoundation with a single observability-noise primitive --- that quantitatively rationalizes the empirical price path. The belief-updating mechanism nests the rational-expectations benchmark as a knife-edge restriction and delivers pre-specified cross-sectional predictions that match the data without additional free parameters.